\section{Leases and employee compensation}

\begin{qfmethod}[Lessee and lessor reporting under IFRS 16 and ASC 842]\label{meth:lessee-lessor-accounting}
\qffield{Statement} A lease conveys control of an identified asset for a period in exchange for consideration. IFRS lessees generally recognize a right-of-use (ROU) asset and lease liability, then depreciation and interest; exemptions exist for short-term and low-value leases. US lessees recognize both finance and operating leases on the balance sheet, but operating leases generally retain a single straight-line expense. Thus the old rule that all operating leases are off balance sheet is unsuitable for the 2026 examples \citep[IFRS~16 paras.~5, 9, 22, 26, 32 and~36]{adopted-accounting-standards-2023}; \citep[ASC~842-20-25-1--25-6, pp.~101--102]{us-leases-2016}.
\qffield{Practice} With three end-year payments of 30, discount rate 6 percent, no residual guarantee, purchase option, initial costs or incentives, the initial liability and ROU asset are 80.1904. First-year interest is 4.8114, principal repayment 25.1886 and closing liability 55.0018. IFRS depreciation is 26.7301, leaving ROU 53.4602 and total first-year expense 31.5415. A US operating lease instead reports expense 30 and ROU amortization 25.1886, leaving ROU 55.0018. Both pay exactly 30; the different expense timing explains the 1.5415 equity difference before tax. IFRS principal and US finance-lease principal are financing flows; US operating-lease payments are generally operating flows. Interest classification follows the applicable cash-flow rules, including the IFRS 18 transition \citep[IFRS~16 para.~50]{adopted-accounting-standards-2023}; \citep[ASC~842-20-45-5, pp.~108--109]{us-leases-2016}.
\qffield{Lessor perspective} An IFRS finance lease transfers substantially all ownership risks/rewards and recognizes a net investment receivable with finance income; an operating lessor retains/depreciates the asset and recognizes rental income. US lessors distinguish sales-type, direct-financing and operating leases; recognition of selling profit, deferred profit and collectibility conditions differs. The lessee's label need not dictate the lessor's classification \citep[IFRS~16 paras.~61--67 and~81]{adopted-accounting-standards-2023}; \citep[ASC~842-30-25-1--25-11, pp.~126--129]{us-leases-2016}.
\qffield{Analysis} Compare lease-adjusted debt and earnings consistently. IFRS/finance presentation raises EBITDA relative to a single operating expense without creating cash. Never add a lease liability again when it is already recognized, and do not deduct both lease payments in cash flow and the same liability in value without a consistent valuation approach.
\end{qfmethod}

\begin{qfmethod}[Defined-contribution and defined-benefit accounting]\label{meth:defined-contribution-benefit-reporting}
\qffield{Statement} In a defined-contribution plan the employer owes agreed contributions for service; employees bear investment and benefit adequacy risk. Expense and any unpaid contribution liability follow the service period. A defined-benefit promise leaves the employer exposed to longevity, salary, discount-rate and investment risk. The balance sheet recognizes the net obligation less plan assets, subject under IFRS to an asset ceiling for a surplus. IFRS profit includes service cost and net interest at the obligation discount rate; remeasurements enter OCI without subsequent recycling. US periodic cost uses an expected asset return and may amortize prior-service and actuarial amounts from accumulated OCI, so identical funded status need not mean identical income \citep[IAS~19 paras.~51, 63--64, 83, 120, 122--127]{adopted-accounting-standards-2023}; \citep[chap.~14, solutions~1--11 and~19--23, PDF pp.~229--231]{statement-analysis-workbook-2020}.
\qffield{Practice} Ignore tax and ceilings; contributions and benefits occur at year-end. Opening obligation 100, assets 80; service cost 8, discount rate 5 percent, benefits 6, employer contribution 10, actual asset return 6 and actuarial obligation loss 4. Obligation closes at $100+8+5+4-6=111$; assets at $80+6+10-6=90$; net liability is 21. IFRS P\&L expense is $8+5-4=9$ and OCI loss $4-(6-4)=2$. Thus $20+9+2-10=21$. US expected asset return of 6 percent gives expense $8+5-4.8=8.2$, assuming no amortization, and OCI loss 2.8; total comprehensive cost is still 11. Cash contribution 10 reduces cash, net liability rises 1 and equity falls 11, so the balance sheet reconciles.
\qffield{Interpretation} Higher expected asset returns can lower US expense without shrinking the benefit promise. For US reporting, service cost belongs with compensation and other components outside operating income if that subtotal is presented; only service cost can enter a constructed asset's cost \citep[ASC~715-20-45-3A, p.~6]{us-pension-2017}. Do not interpret an OCI loss as economically irrelevant merely because net income excludes it.
\end{qfmethod}

\begin{qfmethod}[Cash, share-based and deferred employee compensation]\label{meth:share-based-compensation-reporting}
\qffield{Statement} Salary and short-term cash bonuses consume cash or create payables as service is received. Equity-settled employee awards recognize service cost and an equity credit, generally using grant-date fair value spread over the service period; they consume shareholder value even without a current cash outflow. Cash-settled awards create a liability remeasured through profit to settlement. Pensions defer compensation and introduce the separate risks described in the pension entry \citep[IFRS~2 paras.~7, 11, 15, 19--21 and~30]{adopted-accounting-standards-2023}.
\qffield{Practice} Grant 300 service-vesting shares worth 20 each at grant, vesting after three years. With all expected to vest, total service cost is 6,000 and first-year expense 2,000, reducing retained earnings while increasing contributed equity 2,000; cash is unchanged before tax. If the estimated shares vesting falls to 270 in year two, cumulative expense should be $270(20)(2/3)=3{,}600$, so year-two expense is 1,600 after the first 2,000. A later market price of 30 does not remeasure an unmodified equity-settled grant. A cash-settled award instead remeasures its liability. The examples ignore tax and use simple cliff service vesting.
\qffield{Comparison} IFRS estimates service-condition forfeitures and revises the expected number; US GAAP permits an entity-wide election to estimate or recognize forfeitures as they occur. Market vesting conditions enter grant-date fair value, whereas service and non-market performance conditions affect the awards counted; failing a market target alone does not reverse expense when the other vesting conditions are met. US excess tax benefits/deficiencies generally enter tax expense and operating cash flows under ASU 2016-09, rather than the former equity/financing treatment \citep[Summary pp.~1--3]{us-share-pay-2016}. Both compensation cost and potential dilution matter in assessing performance.
\end{qfmethod}

\begin{qfmethod}[Reading lease, pension, debt and share-pay disclosures]\label{meth:long-term-obligation-disclosures}
\qffield{Statement} Reconcile long-term obligations by amount, maturity, contractual versus estimated payment, security and settlement currency. Debt notes describe rates, maturities, covenants and collateral; lease notes distinguish ROU assets, interest/depreciation or operating lease cost, commitments and cash payments. Pension notes reconcile obligations and assets and identify actuarial assumptions, risk exposures, funding policy and future contributions. Share-pay notes disclose grant types, vesting/settlement terms, option movements, exercise prices, valuation assumptions and expense \citep[IFRS~16 paras.~53 and~89; IAS~19 paras.~135, 140 and~144; IFRS~2 paras.~44--47]{adopted-accounting-standards-2023}; \citep[ASC~842-20-50-4, p.~110]{us-leases-2016}.
\qffield{Practice} The lease example has 90 of undiscounted remaining payments at inception but a recognized liability of 80.1904. A maturity table is therefore not directly comparable to a discounted liability. The pension example has obligation 111 and assets 90, not merely an unexplained net 21; concentrated or illiquid plan assets make the net figure less reassuring. If stock awards reconcile from 1,000 outstanding plus 300 grants minus 80 exercises minus 20 forfeitures to 1,200, confirm those movements against the expense and shares-outstanding schedules rather than assuming 300 grants immediately dilute basic EPS.
\qffield{Interpretation} Obtain covenant definitions rather than substituting a textbook EBITDA measure. Determine whether a waiver existed at the relevant date, whether cash is restricted and whether a surplus pension asset is recoverable. Unrecognized lease commitments, pension sensitivities and equity-award commitments can matter economically even when they are absent from an ordinary debt total. A disclosure gap warrants a range or a request for information, not an invented number.
\end{qfmethod}

\begin{qfmethod}[Forecasting compensation, shares and per-share value]\label{meth:share-pay-shares-forecast}
\qffield{Statement} Forecast share-pay by grant cohort, grant-date value, service period, vesting/forfeiture conditions and settlement form. Reconcile beginning unrecognized compensation plus new grants and modifications less recognized cost and forfeitures to the remaining balance. Forecast actual shares separately from expense: vesting, option exercise, repurchases and tax withholding change share counts on their own dates. Basic EPS uses weighted-average shares; potential dilution follows the separate diluted-EPS rules \citep[IFRS~2 paras.~15, 19--21, 44--47; IAS~33 paras.~19, 31 and~45]{adopted-accounting-standards-2023}.
\qffield{Practice} Existing unrecognized service cost 6,000 has two years remaining; a new start-year grant worth 6,000 vests over three years. With straight-line service and no forfeitures, next-year expense is $3{,}000+2{,}000=5{,}000$. Beginning shares 1,000, a start-year repurchase of 20 and vesting of 100 new shares halfway through the year give closing shares 1,080 and weighted-average shares 1,030. At market price 30 the buyback consumes 600 cash; it does not erase the 5,000 compensation expense. Model a 10 percent forfeiture change or a higher grant value as a sensitivity rather than fixing SBC as an unexplained balancing percentage.
\qffield{Valuation} A forecast that treats recurring SBC as an economic operating cost can retain that cost in free cash flow, while valuing existing awards and share dilution consistently. A model that adds SBC back as noncash must explicitly charge the economic cost of future awards through dilution or an equivalent cash replacement cost. Do not both deduct the same award's full fair value and count its full dilution without reconciling the overlap. Nor should every noncash expense be called free: employees received value and existing owners bear it. Distinguish the equity value for all claimants from value per current common share.
\end{qfmethod}

\begin{qfmethod}[Forecasting pension balances, cash funding and valuation]\label{meth:pension-model-valuation}
\qffield{Statement} Forecast the gross benefit obligation and plan assets before the net amount. Obligation changes include service, interest, benefits, actuarial remeasurement and plan amendments; assets change with actual return, employer contributions and benefits paid. Accounting expense is not the cash funding requirement. Expected returns, discount curves, salary/longevity assumptions and contribution policy should be explicit and internally consistent \citep[IAS~19 paras.~83, 120, 123, 127, 140 and~144]{adopted-accounting-standards-2023}; \citep[chap.~14, solutions~8--15 and~23, PDF pp.~230--232]{statement-analysis-workbook-2020}.
\qffield{Practice} Starting from obligation 111 and assets 90, assume next-year service 8.4, discount rate 5 percent, actual asset return 4.5, contributions 12 and benefits 7, all year-end flows, with no remeasurement. Obligation closes at $111+8.4+5.55-7=117.95$; assets at $90+4.5+12-7=99.5$; deficit is 18.45. IFRS expense is $8.4+0.05(21)=9.45$, so $21+9.45-12=18.45$. The 12 contribution is an operating cash use in the stated model, not 9.45. A fall in the discount rate generally raises the obligation; a negative asset return reduces plan assets even if US expected-return income remains positive.
\qffield{Valuation} For an operating enterprise valuation, keep recurring service cost in operations, separate pension financing/remeasurements, and consider the underfunded claim in the equity bridge. With deficit 21, an assumed fully usable 25 percent future tax benefit would give 15.75 after tax; actual recovery and timing require evidence. Alternatively model all future deficit-repair contributions directly. Subtracting the deficit and also deducting the same repair contributions from forecast free cash flow double counts that burden. A reported surplus should not be added at face value without checking refunds, restrictions and IFRS's asset ceiling.
\end{qfmethod}
